\documentclass{book}
\usepackage[T2A]{fontenc}
\usepackage[a5paper]{geometry}
\usepackage[utf8]{inputenc}
\usepackage[english, russian]{babel}
\usepackage{fancyhdr}
\usepackage{amsthm}
\usepackage{breqn}
\usepackage{amsmath}
\usepackage{setspace}
\usepackage{amssymb}
\usepackage{geometry}
\setlength{\parskip}{1em}
\geometry{verbose,a5paper,tmargin=1cm,lmargin=2.5cm,rmargin=2.5cm}
\linespread{0.9}
\setcounter{page}{229}
\cfoot{\overline{\quad\quad\textsl{\thepage}\quad\quad}}
\renewcommand{\headrulewidth}{0pt}
\pagestyle{fancy}

\begin{document}
\noindent выполняется неравенство
\begin{equation}
|f(x')-f(x)|<\varepsilon
\tag{6.14}
\end{equation}
\par Если число \begin{math} \delta \end{math} можно выбрать не зависящим от точки $x$, так, чтобы при выполнении условия (6.13) выполнялось условие (6.14), то функция $f$ называется равномерно непрерывной. Сформулируем определение этого важного понятия более подробно.\\
\textsf{\textbf{\small{О п р е д е л е н и е 3.}}}
\textit{Функция $f$, заданная на отрезке}[$a, b$], \textit{называется равномерно непрерывной на нем, если для любого} 
$\varepsilon>0$ \textit{существует такое} $\delta>0$, \textit{что для любых двух точек} 
\begin{math} x \in[$a,b$]\end{math}
\textit{ и} \begin{math}x' \in[$a,b$]\end{math} \textit{таких, что} \begin{math}|x'-x| < \delta \end{math}\textit{ выполняется неравенство} \begin{math}|f(x')-f(x)|<\varepsilon \end{math}.
\par В символической записи определение непрерывности функции на отрезке выглядит следующим образом:
\begin{equation} \notag
\forall \: x \: \forall \: \varepsilon>0 \: \exists \: \delta>0 \forall \: x', |x'-x|<\delta: |f(x')-f(x)|<\varepsilon,
\end{equation} а определение равномерной непрерывности так:
\begin{equation} \tag{6.15}
\forall \: \varepsilon>0 \: \exists \: \delta>0 \: \forall \: x, x', |x'-x|<\delta: |f(x')-f(x)|<\varepsilon.
\end{equation}
Здесь точки $x$ и $x'$ принадлежат отрезку, на котором рассматривается функция $f$.
\par Ясно, что всякая равномерно непрерывная на отрезке функция непрерывна на нем: если в определении равномерной непрерывности зафиксировать точку $x$, то получится определение непрерывности в этой точке.
\par\textsf{\textbf{\small{П Р И М Е Р Ы. 1.}}}
Функция \begin{math}f(x)=x\end{math} равномерно непрерывна на всей числовой оси, так как, если задано 
\begin{math}\varepsilon > 0\end{math}, достаточно взять \begin{math}\delta = \varepsilon \end{math}; тогда если 
\begin{math}|x-x'| < \delta \end{math}, то, в силу равенств 
\begin{math} f(x) = x \end{math}, \begin{math} f(x') = x' \end{math}, получим
\begin{math} |f(x)-f(x')| < \varepsilon \end{math}.
\par \textsf{\textbf{\small{2.}}} Функция \begin{math}f(x) = sin(1/x), x \neq 0\end{math},
будучи непрерывной на своей области определения, т.е. на числовой оси, из которой удалена точка 
\begin{math}x = 0\end{math}, не будет на ней равномерно непрерывна.
\par В самом деле, если взять, например, \begin{math}\varepsilon = 1\end{math}, то при любом сколь угодно малом \begin{math}\delta > 0\end{math} найдутся точки $x$ и $x'$, например точки вида
\begin{math}x = \frac{1}{(\pi/2 + 2 \pi n)}\end{math} и \begin{math}x' = \frac{1}{(3\pi/2 + 2 \pi n)}\end{math}
(n - достаточно большое натуральное число) такие, что \begin{math}|x - x'| < \delta \end{math}, а вместе с тем
\begin{math}|f(x) - f(x')| > 1\end{math}.
\par \textsf{\textbf{\small{3.}}} Функция \begin{math}f(x) = x^2\end{math} не равномерно непрерывна на всей числовой оси \R. Это следует из того, что для любого фиксированного \begin{math}h \neq 0\end{math} имеет место равенство
\begin{equation} \notag
\lim_{x\to \infty} 	[f(x + h) - f(x)] = \lim_{x\to \infty} [(x + h)^2 - x^2] = \lim_{x\to \infty} [2xh + h^2] = \infty.
\end{equation}
Поэтому если задано \begin{math}\varepsilon > 0\end{math}, то, каково бы ни было \begin{math}\delta > 0\end{math}, зафиксировав \begin{math}h \neq 0, |h| < \delta \end{math}, можно так выбрать $x$, что будет выполняться неравенство
\begin{math}|f(x+h) - f(x)| > \varepsilon \end{math}.
\par\textsf{\textbf{\small{Т Е О Р Е М А 5 (Кантора).}}} Функция, непрерывная на отрезке, равномерно непрерывна на нём. \medskip
\\ \textsf{\small{Д О К А З А Т Е Л Ь С Т В О.}} Докажем теорему от противного. Допустим, что на некотором отрезке
[$a,b$] существует непрерывная, однако не равномерно непрерывная на нем функция $f$. Это означает (см.(6.15)), что существует такое \begin{math}\varepsilon_{0} > 0\end{math}, что для любого \begin{math}\delta > 0\end{math} найдутся такие точки \begin{math}x \in [a, b]\end{math} и \begin{math}x' \in [a, b]\end{math}, что \begin{math}|x'-x| < \delta \end{math}, но \begin{math}|f(x') - f(x) \geqslant \varepsilon_{0} \end{math}. В частности, для 
\begin{math}\delta = 1 / n\end{math} найдутся такие точки, обозначим их $x_{n}$ и $x'_{n}$, что
\begin{equation} \tag{6.16}
|x'_{n} - x_{n}| < \frac{1}{n},
\end{equation} но
\begin{equation} \tag{6.17}
|f(x'_{n}) - f(x_{n})| \geqslant \varepsilon_{0}.
\end{equation}
\indent Из последовательности точек \begin{math}\{x_{n}\}\end{math} в силу свойства компактности (см. теорему 4 в п.4.6) можно выделить сходящуюся подпоследовательность \begin{math}\{x_{n_{k}}\}\end{math}. Обозначим ее предел $x_{0}$:
\begin{equation} \tag{6.18}
\lim_{k\to \infty} x_{n_{k}} = x_{0}.
\end{equation}
Поскольку \begin{math}a\leqslant x_{n_{k}} \leqslant b, k = 1,2, ...\end{math}, то 
\begin{math}a\leqslant x_{0} \leqslant b\end{math} (см. п.4.3). Функция $f$ непрерывна в точке $x_{0}$, поэтому
\begin{equation} \tag{6.19}
\lim_{k\to \infty} f(x_{n_{k}}) \underset{(6.18)} = f(x_{0}).
\end{equation}
Подпоследовательность $\{x'_{n_{k}}\}$ последовательности $\{x'_{n}\}$ также сходится к точке $x_{0}$, ибо при $k \to \infty$.
\begin{equation} \notag
|x'_{n_{k}} - x_{0}| \geqslant |x'_{n_{k}} - x_{n_{k}}| + |x_{n_{k}} - x_{0}| \underset{(6.16)} < \frac{1}{n_{k}} + |x_{n_{k}} - x_{0}| \underset{(6.18)} \to 0.
\end{equation}

\end{document}